\documentclass{article}

\usepackage[T1]{fontenc}
\usepackage{lmodern}
\usepackage{amsmath}
\usepackage{amssymb}
\usepackage{amsthm}
\usepackage{mathtools}
\usepackage[hidelinks]{hyperref}

\newtheorem{proposition}{Proposition}

\DeclareMathOperator{\rank}{rank}
\DeclareMathOperator{\supp}{supp}
\newcommand{\Ftwo}{\mathbb{F}_2}

\title{An explicit rectangular CSS model on 722 qubits: Commutation, independent boundary checks, and quotient dimension}
\author{}
\date{Frozen document-review checkpoint, 11 November 2026}

\begin{document}

\maketitle

\textbf{Claim grades.} THEOREM means refereed complete; PROPOSITION means a claim
with a named gap; CONJECTURE means a mathematically unproved claim; FALSE means a refuted
claim. Every derived claim in this draft is PROPOSITION: the named gap is independent
review\slash formalization pending. The arguments below are offered as complete proofs of those
claims, not as referee reports. No claim is graded THEOREM. This frozen checkpoint is ready
for three independent document-only referees. Their reviews are not yet completed. A final
PROOF may contain only claims that have passed full independent review and formalization;
this file is not that final PROOF.

\section{Main result and explicit finite model}\label{sec:model}

All vector and polynomial arithmetic is over \(\Ftwo\). Let
\begin{equation}
  Q = \{0, \dots, 18\} \times \{0, \dots, 18\}.
\end{equation}
There are two qubits per cell, denoted \(A(i,j)\) and \(B(i,j)\), with indices
\(19i + j\) and \(361 + 19i + j\). Thus the physical vector space is
\(V = F^{722}\), where here and below \(F = \Ftwo\). Encode a vector by a pair
of polynomials \((u, v)\) with exponents in \(Q\), one for each block. The
coefficients are its physical bits.

Finite Laurent polynomials allow integer exponents. Put
\begin{equation}
  f = x + x^2 + y^2, \quad g = 1 + x^2 y + x^2 y^2,
\end{equation}
and write \(\overline{p}(x,y) = p(x^{-1}, y^{-1})\). Let \(P\) delete every
monomial whose exponent pair is outside \(Q\), in each component separately.
Define two binary matrices by their rows, in lexicographic anchor order:
\begin{align}
  H_X[(a,b), {:}] &= P\bigl(x^a y^b (f, g)\bigr),
  & -2 &\le a \le 18, & 0 &\le b \le 16; \\
  H_Z[(c,d), {:}] &= P\bigl(x^c y^d (\overline{g}, \overline{f})\bigr),
  & 2 &\le c \le 18, & 0 &\le d \le 20.
\end{align}
Rank means dimension of the row span. For a subspace \(S\), write
\(S^{\perp} = \{w \in V : s \cdot w = 0 \text{ for every } s \in S\}\), where
the dot product is the sum of coordinate products in \(F\). Both matrices have
\(357\) rows and \(722\) columns. Equivalently a row bit is 1 exactly when the
indicated shifted support contains that qubit and the qubit lies in \(Q\). No
external file or construction algorithm is part of this definition.

\begin{proposition}[main result; independent review\slash formalization pending]\label{prop:main}
  For these explicit matrices,
  \begin{equation}
    H_X H_Z^T = 0, \quad \rank(H_X) = \rank(H_Z) = 357.
  \end{equation}
  Writing \(S_X\) and \(S_Z\) for their row spaces, the quotient
  \(S_Z^{\perp}/S_X\) has dimension \(8\). Here \(k = 8\) denotes only this binary
  quotient dimension. A Hilbert-space interpretation is outside the scope of this review
  checkpoint.
\end{proposition}

The proof follows in Sections~\ref{sec:polynomial}--\ref{sec:quotient}. All
auxiliary algebra and dimension facts needed for it are proved there.

\section{Polynomial facts used in the proof}\label{sec:polynomial}

\begin{proposition}[elementary algebra; independent review\slash formalization pending]%
  \label{prop:algebra}
  The rings \(F[z]\) and \(F[z, z^{-1}]\) have no zero divisors. In either
  ring, extreme exponents of a nonzero product add. Polynomial equality is
  coefficientwise, and multiplication is finite coefficient convolution.
\end{proposition}

\begin{proof}
  A polynomial is by definition a finite coefficient function on nonnegative
  integer exponents; a Laurent polynomial uses integer exponents. Addition is
  coefficientwise and multiplication is convolution. If \(p\) and \(q\) have
  largest nonzero exponents \(r\) and \(s\), the coefficient at \(r + s\) in
  \(pq\) is the product of their two leading coefficients. The only pair of
  exponents attaining this sum is \((r, s)\); its coefficient is 1 in \(F\).
  Thus the product is nonzero, with largest exponent \(r + s\). The same
  argument applies to smallest exponents. It also proves the statement over any coefficient ring
  without zero divisors, using a nonzero product of extreme coefficients.
  Iterating this observation allows polynomials in \(x\) with coefficients in
  \(F[y]\), and vice versa. In particular a nonzero polynomial times an
  unknown polynomial can vanish only if the unknown vanishes. No unique
  factorization or assertion about polynomial ideals is needed.
\end{proof}

\section{Commutation without a corner assumption}\label{sec:commutation}

Consider first the untruncated rows on the integer lattice. The overlap parity
of the A components at anchors \((a,b)\) and \((c,d)\) is the coefficient of
\(x^{c-a} y^{d-b}\) in \(fg\): an intersection is a pair of support exponents
\(s\) in \(f\), \(t\) in \(g\), with \((a,b) + s = (c,d) - t\). The B overlap
is the same coefficient in \(gf\). Their sum is zero because convolution
commutes and \(1 + 1 = 0\) in \(F\).

Every full X row has all its second coordinates between 0 and 18, because
\(0 \le b \le 16\) and all second support exponents of \(f, g\) lie between 0
and 2. Every full Z row has all its first coordinates between 0 and 18,
because \(2 \le c \le 18\) and all first support exponents of
\(\overline{g}, \overline{f}\) lie between \(-2\) and 0. Consequently any
intersection of a full X row and a full Z row already belongs to \(Q\).
Truncating both rows removes no intersection. Their finite overlap is
therefore zero. This proves \(H_X H_Z^T = 0\) for every allowed pair of
anchors; no corner completion is assumed.

\section{Independence by two one-dimensional recurrences}\label{sec:independence}

\subsection{X checks}\label{sec:x-checks}

Let a linear combination of X rows have anchor polynomial
\begin{equation}
  h(x,y) = \sum_{i=-2}^{18} h_i(y) x^i, \quad
  \supp(h_i) \subseteq \{0, \dots, 16\}.
\end{equation}
Suppose its vector \(P(fh, gh)\) is zero. Neither product has a second
coordinate outside 0 through 18. Thus for each \(i = 0, \dots, 18\) all
coefficients in the first-coordinate slice vanish, as polynomials in \(F[y]\):
\begin{equation}
  h_{i-2} + h_{i-1} + y^2 h_i = 0, \quad (y + y^2) h_{i-2} + h_i = 0.
\end{equation}
Write \(t = y + y^2\) and \(b = 1 + y^2 t = 1 + y^3 + y^4\). Substitution gives
\begin{equation}
  h_{i-1} = b h_{i-2}.
\end{equation}
At \(i = 0, 1\) this says \(h_{-1} = b h_{-2}\) and \(h_0 = b^2 h_{-2}\),
whereas the second original equation gives \(h_0 = t h_{-2}\). Hence
\begin{equation}
  (b^2 + t) h_{-2} = (1 + y + y^2 + y^6 + y^8) h_{-2} = 0.
\end{equation}
The displayed multiplier is nonzero. Proposition~\ref{prop:algebra} gives
\(h_{-2} = 0\), then \(h_{-1} = 0\). Induction in \(h_i = t h_{i-2}\) gives
\(h_i = 0\) for all \(0 \le i \le 18\). Every anchor coefficient is zero.
Therefore the 357 X rows are independent.

\subsection{Z checks}\label{sec:z-checks}

Reflect physical coordinates by \((i,j) \mapsto (18 - i, 18 - j)\) and
exchange the two blocks. This is a permutation, so a vector vanishes before it
exactly when it vanishes after it. It sends a Z row at \((c,d)\) to
\(P\bigl(x^{18-c} y^{18-d} (f, g)\bigr)\). Thus a vanishing Z combination
becomes \(P(fh, gh) = 0\) with
\begin{equation}
  h(x,y) = \sum_{j=-2}^{18} h_j(x) y^j, \quad
  \supp(h_j) \subseteq \{0, \dots, 16\}.
\end{equation}
This permutation is used to prove independence, not asserted to preserve the
code's check sets. No first-coordinate truncation occurs in these products.
Set \(a = x + x^2\). For \(j = 0, \dots, 18\) vanishing gives
\begin{equation}
  h_{j-2} + a h_j = 0, \quad h_j + x^2 h_{j-1} + x^2 h_{j-2} = 0.
\end{equation}
With \(b = 1 + x^2 a = 1 + x^3 + x^4\), substitution yields
\begin{equation}
  x^2 h_{j-1} = b h_j.
\end{equation}
The equations at \(j = 0, 1\) imply \(x^4 h_{-1} = b^2 h_1\). The first
original equation at \(j = 1\) gives \(h_{-1} = a h_1\). Therefore
\begin{equation}
  (b^2 + x^4 a) h_1 = (1 + x^5 + x^8) h_1 = 0.
\end{equation}
Proposition~\ref{prop:algebra} forces \(h_1 = 0\). The relation at \(j = 1\) then
gives \(h_0 = 0\); the first equations give \(h_{-2} = h_{-1} = 0\). For
\(j \ge 2\), induction in \(a h_j = h_{j-2}\) gives \(h_j = 0\), because \(a\)
is nonzero. Hence all 357 Z rows are independent. This completes the rank part
of Proposition~\ref{prop:main}.

\section{Quotient dimension}\label{sec:quotient}

We supply the finite linear algebra explicitly. The span of an independent
list of \(r\) vectors has dimension \(r\): unique coordinates identify it with
\(F^r\). It has \(2^r\) elements, so two bases of the same binary vector space
necessarily have the same size. This also justifies the dimension convention.
An independent list in a finite vector space extends to a basis by repeatedly
adjoining a vector outside its span; the process stops because the space has
finitely many vectors. For a subspace \(U\) of dimension \(r\) in \(F^n\),
extend a basis of \(U\) to a basis of \(F^n\). The images of the added vectors
form a basis of the quotient, proving \(\dim(F^n/U) = n - r\).

The dot product on \(F^n\) is nondegenerate: a nonzero vector has a coordinate
1 and pairs nontrivially with that coordinate's unit vector. If rows
\(z_1, \dots, z_r\) are independent, the functionals \(v \mapsto z_i \cdot v\)
are independent, since any linear relation between them, evaluated on all unit
vectors, is the same relation between the rows. Their combined map
\(T \colon F^n \to F^r\) is onto. Indeed, if its image were a proper subspace,
extend a basis of that image to a basis of \(F^r\) and choose a nonzero linear
functional zero on the image, by prescribing its values on this basis. In
standard coordinates that functional is a nonzero linear combination of the
\(r\) coordinates, contradicting their independence.

Choose preimages of a basis of \(F^r\). Together with a basis of \(\ker T\)
they form a basis of \(F^n\): applying \(T\) proves independence, and
subtracting the appropriate preimages proves spanning. Thus
\(\dim(\ker T) = n - r\). Apply this to \(H_Z\):
\(\dim S_Z^{\perp} = 722 - 357 = 365\). Commutation proved
\(S_X \subseteq S_Z^{\perp}\); the quotient basis argument gives
\(\dim(S_Z^{\perp}/S_X) = 365 - 357 = 8\).

This completes Proposition~\ref{prop:main}. \qed

\end{document}
